\section*{引言}
我们当前正在见证理解与研究量子多体物理方式的一场范式转变（paradigm shift）。在过去一个世纪的大部分时间里，\emph{量子多体问题}（quantum many-body problem）一直是理论物理与量子化学的核心问题，而与之相伴的指数墙（exponential wall）阻碍了描述强相互作用物质相（phases of matter）的进展。然而在过去的二十年中，基于\emph{纠缠}（entanglement）理论来模拟与分类这些强关联系统（strongly correlated systems）取得了显著突破，其典范代表正是\emph{张量网络}（tensor network）的概念。

从本质上讲，张量网络通过提供多体波函数的\emph{全息}（holographic）描述，绕过了指数级的量子多体墙。张量网络并非在指数级巨大的希尔伯特空间（Hilbert space）中描述波函数，而是对系统不同部分之间纠缠的路由方式进行建模。张量网络之所以取得成功，其根本原因在于：计算低能波函数任意可观测量所需的全部信息，都可以压缩至局域张量（local tensor）中，其中所谓的\emph{纠缠自由度}（entanglement degrees of freedom）构成了\emph{虚}（virtual）希尔伯特空间。能够做到这一点的根源在于纠缠熵\emph{面积律}（area law）的存在：张量网络正是对所有满足此类面积律的态构成的流形（manifold）的简洁而精确的表示。通过将定义在指数级巨大物理希尔伯特空间中的哈密顿量（Hamiltonian）投影到该低维流形上，构建用于模拟大规模量子多体系统的数值算法成为了可能。DMRG（密度矩阵重整化群，\emph{Density Matrix Renormalisation Group}）、MPS（矩阵乘积态，\emph{Matrix Product State}）与 PEPS（投影纠缠对态，\emph{Projected Entangled-Pair State}）算法正在彻底改变我们模拟原子物理前沿实验、当今量子计算机以及量子多体系统低能物理的方式。

不仅如此：自 Landau 与 Anderson 的开创性见解以来，人们便认识到多体物理的核心在于\emph{对称性}（symmetries）与\emph{对称性破缺}（symmetry breaking）。事实证明，整个系统的全局、局域及拓扑对称性，都反映在张量网络局域张量的对称性之中，而这些局域张量对整个多体波函数提供了极其有效的压缩。因此，对有能隙（gapped）\emph{物质相}、特别是\emph{拓扑相}（topological phases）的分类，便归结为对张量在一组相关对称性下不可约变换方式的分类问题。由于这些对称性作用在虚（纠缠）自由度上，并无理由要求它们必须构成一个群，从而涌现出可以用\emph{融合范畴}（fusion categories）来描述的更为一般的结合代数结构（associative structures）。事实上，张量网络提供了一种显式方法，通过 MPO（矩阵乘积算符，\emph{Matrix Product Operator}）代数及其高维类似物 PEPO（投影纠缠对算符，\emph{projected entangled-pair operators}）来表示这些\emph{广义对称性}（generalised symmetries）以及它们之间的对偶性（dualities）。在数学物理、凝聚态物理、统计物理以及高能物理中，有若干最活跃的研究方向正致力于探索这一深刻见解所带来的推论：强关联与拓扑物质相的全息描述，是唯一能够揭示其内在结构的描述途径。

\bigskip\noindent 本讲义基于 Frank Verstraete 于 2025 年 8 月在 Les Houches 物理暑期学校上以“\emph{精确可解性与量子信息}”（\emph{Exact Solvability and Quantum Information}）为题所作的五次讲座，涵盖了张量网络若干基础的计算与理论成果。本讲义并不旨在对这些主题进行详尽无遗的综述，而是作为深入探索的引导；读者亦可参考现存大量优秀的综述与讲义文献：~\cite{Verstraete2008,Schollwoeck2010,Bridgeman2016,Vanderstraeten2019,Cirac2020,Xiang2023,Banuls2022,Verstraete2023,Chen2025}。